% ===================================================================
%  TAULA N.º 12 — L'Estació. — Els Ferrocarrils. — Els Vaixells.
%  Traduction 2026 d'apres texto/io, controlee sur texto/fr et
%  texto/es. Consigne : voir texto/ca/10-taula-01.tex.
% ===================================================================

%%K t12-apar-1 apar
\VUsaut{4.0mm}
\VUtitre{15.0pt}{60}{TAULA N.º 12}
\VUsaut{2.4mm}
\VUpk{11.4pt}{40}{L'Estació. --- Els Ferrocarrils. --- Els Vaixells.}
\VUsaut{3.4mm}
\textsc{Nota.} --- \textit{Els horaris usats a cada país són diferents, per la qual
cosa hem pres com a exemple les hores del} \VUgras{quadre francès}.

%%K t12-01-1 p
1. --- Acabo de fer un viatge per Alemanya. Abans d'emprendre'l havia comprat una
\VUgras{guia d'horaris} \textsuperscript{(15)} per saber l'hora de sortida dels
trens. Havent comprovat que el tren més còmode sortia de París a les nou de la
nit i arribava a Estrasburg a la una i tretze, vaig preparar el meu viatge i,
l'endemà, em vaig fer portar a l'estació.

%%K t12-02-1 p
2. --- Quan vaig entrar a la gran sala de sortides, darrere d'un vell
\VUgras{rector} \textsuperscript{(2)} de poble que portava a la mà la seva
\VUgras{bossa de viatge} \textsuperscript{(3)}, ja havien arribat molts
\VUgras{viatgers} \textsuperscript{(1)}.

%%K t12-02-2 p
Com que volia enviar un \VUgras{telegrama} \textsuperscript{(5)}, vaig fer que un
\VUgras{revisor} \textsuperscript{(6)} m'indiqués l'\VUgras{oficina de telègrafs}
\textsuperscript{(4)}.

%%K t12-03-1 p
3. --- Abans de passar a la \VUgras{sala d'espera} \textsuperscript{(7)} vaig anar a
treure un \VUgras{bitllet} \textsuperscript{(9)} d'anada i tornada.

%%K t12-04-1 p
4. --- No vaig esperar gaire a la \VUgras{taquilla} \textsuperscript{(8)}, perquè hi
havia poca gent. Un \VUgras{soldat} \textsuperscript{(10)} que sortia de permís em
va cedir amablement el seu torn, de manera que vaig tenir temps de demanar
algunes dades al \VUgras{cap d'estació} \textsuperscript{(13)}, que parlava amb el
\VUgras{comissari} \textsuperscript{(12)} de vigilància i amb el
\VUgras{gendarme} \textsuperscript{(11)} de servei.

%%K t12-tit-1 sub
\VUpk{10.2pt}{40}{Facturar l'equipatge.}
\VUsaut{3.0mm}

%%K t12-05-1 p
5. --- Després de comprar un diari al \VUgras{quiosc de llibres}
\textsuperscript{(14)}, vaig fer pesar el meu \VUgras{equipatge}
\textsuperscript{(17)}, que un \VUgras{mosso} \textsuperscript{(16)} acabava de
portar en un \VUgras{carretó} \textsuperscript{(19)}; l'\VUgras{empleat}
\textsuperscript{(22)} encarregat de la \VUgras{bàscula} \textsuperscript{(21)} em
va dir que el meu bagul pesava trenta-cinc quilos, la qual cosa feia un
\VUgras{excés} de cinc quilos, pel qual havia de pagar un suplement a la
\VUgras{finestreta d'equipatges} \textsuperscript{(23)}.

%%K t12-06-1 p
6. --- Com que anava amb avançament, vaig tenir ocasió de mirar el moviment dels
viatgers per l'estació. Uns deixaven o recollien petits embalums a la
\VUgras{consigna} \textsuperscript{(25)}; d'altres consultaven els \VUgras{horaris}
\textsuperscript{(26)}. Els viatgers que arribaven s'afanyaven cap a la
\VUgras{sortida} \textsuperscript{(32)} amb la seva \VUgras{maleta}
\textsuperscript{(24)} a la mà. Alguns esperaven al \VUgras{lliurament
d'equipatges} \textsuperscript{(27)} i presentaven el seu \VUgras{resguard}
\textsuperscript{(29)} per retirar els seus baguls, \VUgras{caixes}
\textsuperscript{(28)} o \VUgras{cistelles} \textsuperscript{(30)}.

%%K t12-07-1 p
7. --- Els \VUgras{empleats de consums} \textsuperscript{(33)} registraven amb cura
els embalums per veure si hi havia res a declarar.

%%K t12-08-1 p
8. --- Alguns viatgers es dirigien a la \VUgras{fonda} \textsuperscript{(34)}, on un
\VUgras{cambrer} \textsuperscript{(35)} s'afanyava a servir-los. Havien de
donar-se pressa, perquè el \VUgras{tren} \textsuperscript{(36)}, que era un exprés,
no es quedava gaire temps a l'estació.

%%K t12-09-1 p
9. --- Vaig passar després a l'andana de sortida i vaig buscar un \VUgras{vagó}
\textsuperscript{(37)} directe per no haver de canviar durant el viatge. Vaig
pujar a un cotxe de passadís que té un \VUgras{departament}
\textsuperscript{(38)} per a fumadors i un altre reservat a «senyores soles».

%%K t12-tit-2 sub
\VUpk{10.2pt}{40}{Sortida de nit.}
\VUsaut{3.0mm}

%%K t12-10-1 p
10. --- Després de pujar l'\VUgras{estrep} \textsuperscript{(40)}, tancar la
\VUgras{porteta} \textsuperscript{(39)}, apujar la \VUgras{cortineta}
\textsuperscript{(41)} i col·locar els meus embalums petits a la \VUgras{reixeta}
\textsuperscript{(43)}, em vaig asseure al \VUgras{seient}
\textsuperscript{(42)}. En veure el \VUgras{fanaler} \textsuperscript{(44)}
encendre els llums, vaig pensar que tenia una nit per passar al vagó i vaig
cridar l'empleat encarregat del lloguer dels \VUgras{coixins}
\textsuperscript{(45)} per demanar-ne un.

%%K t12-11-1 p
11. --- El conductor del tren va venir a comprovar els bitllets. Li vaig preguntar
per què no sortíem, ja que era l'hora. Em va contestar que el \VUgras{cap de tren}
\textsuperscript{(46)} estava ocupat fent col·locar unes bicicletes al
\VUgras{furgó} \textsuperscript{(47)}. A més, encara no s'havien canviat tots els
\VUgras{escalfapeus} \textsuperscript{(48)}.

%%K t12-12-1 p
12. --- Però no trigaríem a sortir, perquè el \VUgras{fogoner}
\textsuperscript{(50)}, damunt el seu \VUgras{tènder} \textsuperscript{(49)},
agafava carbó per revifar el foc de la \VUgras{locomotora}
\textsuperscript{(54)}, i \VUgras{el maquinista} \textsuperscript{(51)} només
esperava el senyal del \VUgras{cap adjunt d'estació} \textsuperscript{(52)}, que
era a l'\VUgras{andana} \textsuperscript{(53)}.

%%K t12-13-1 p
13. --- La \VUgras{caldera} \textsuperscript{(56)} de la locomotora rugia
sordament; la \VUgras{xemeneia} \textsuperscript{(55)} vomitava torrents de fum
barrejat amb \VUgras{vapor} \textsuperscript{(60)}. Va sonar un \VUgras{xiulet}
\textsuperscript{(59)} agut i els \VUgras{pistons} \textsuperscript{(58)} es van
posar en moviment, accionant les \VUgras{rodes} \textsuperscript{(57)} de la
pesada màquina.

%%K t12-tit-3 sub
\VUsaut{3.0mm}
\VUpk{10.2pt}{40}{En marxa!}
\VUsaut{3.0mm}

%%K t12-14-1 p
14. --- La velocitat del tren, que corria damunt els \VUgras{rails}
\textsuperscript{(64)}, es va accelerar; les \VUgras{andanes}
\textsuperscript{(65)} es van acabar; vam passar els \VUgras{comuns}
\textsuperscript{(66)} i també una filera de vagons apartats en una altra
\VUgras{via} \textsuperscript{(63)}: \VUgras{cotxes llit} \textsuperscript{(67)},
\VUgras{cotxe restaurant} \textsuperscript{(68)}, \VUgras{cotxe correu}
\textsuperscript{(69)}.

%%K t12-noto-1 noto
\VUnotes{143.2mm}{%
\textit{(*) Nota.} --- \textit{S'entén que el viatge es descriu tal com era al
llarg de la frontera d'abans de la guerra.}%
}

%%K t12-15-1 p
15. --- Després va passar davant dels nostres ulls l'\VUgras{estació de
mercaderies} \textsuperscript{(71)}. Sota els coberts, uns empleats carregaven les
\VUgras{mercaderies} \textsuperscript{(72)} enviades a petita velocitat. Uns
\VUgras{mossos de maniobres} \textsuperscript{(76)}, vora el \VUgras{dipòsit
d'aigua} \textsuperscript{(73)}, maniobraven \VUgras{vagons de bestiar}
\textsuperscript{(75)} i \VUgras{vagons plataforma} \textsuperscript{(79)} per
formar un \VUgras{tren de mercaderies} \textsuperscript{(80)}.

%%K t12-16-1 p
16. --- Vam passar l'última \VUgras{agulla} \textsuperscript{(78)}, vora la qual
era el guardaagulles; vam deixar enrere el \VUgras{disc de senyals}
\textsuperscript{(86)} i l'estació es va perdre en la llunyania.

%%K t12-17-1 p
17. --- Aviat vam entrar al \VUgras{pont de ferro} \textsuperscript{(74)} que creua
el \VUgras{riu} \textsuperscript{(81)}. Passat el pont, vam seguir el curs del
riu, pel qual potents \VUgras{remolcadors} \textsuperscript{(82)} arrossegaven
pesades \VUgras{gavarres} \textsuperscript{(83)}. Vam veure a més un
\VUgras{vaixell de passatgers} \textsuperscript{(84)} en el moment en què anava a
atracar al \VUgras{pontó} \textsuperscript{(85)}.

%%K t12-18-1 p
18. --- Després de passar de llarg diverses estacions, vam travessar alguns
\VUgras{túnels} \textsuperscript{(87)} i vam deixar darrere nostre innombrables
\VUgras{casetes} \textsuperscript{(88)} de \VUgras{guardabarreres}. Bressolat pel
sotragueig del tren, em vaig quedar profundament adormit.

%%K t12-tit-4 sub
\VUsaut{3.0mm}
\VUpk{10.2pt}{40}{L'estació de Deutsch-Avricourt.}
\VUsaut{3.0mm}

%%K t12-19-1 p
19. --- Em va despertar de sobte la veu d'un empleat que cridava:
«Deutsch-Avricourt! \textsuperscript{(*)} Baixin tots!». Era la inspecció de
duanes i totes les enutjoses formalitats de la frontera. Vaig haver de posar el
meu rellotge en hora amb el \VUgras{rellotge} \textsuperscript{(89)} de l'estació,
que anava cinquanta-cinc minuts avançat; amb prou feines despert, distingia amb
dificultat l'\VUgras{agulla gran} \textsuperscript{(90)} de la \VUgras{petita}
\textsuperscript{(91)}; l'\VUgras{esfera} \textsuperscript{(92)} mateixa estava del
tot confusa per la boira del matí.

%%K t12-20-1 p
20. --- Mentre els duaners feien la seva inspecció, vaig desxifrar badallant els
\VUgras{cartells} \textsuperscript{(93)} que decoren les parets de l'estació. Vaig
esmorzar per fer temps, vaig consultar el mapa de la \VUgras{xarxa}
\textsuperscript{(94)} alemanya per veure quina distància em separava encara
d'Estrasburg, i vaig tornar a pujar al vagó, reconfortat per l'esperança
d'arribar aviat al terme del meu viatge.
\VUsaut{6.0mm}
\VUfilet{22mm}
